\section{Problem Statement}

\subsection{Knowledge Graph Error}

We define a knowledge graph \( \mathcal{G} \) as a set of triples \( (h, r, t) \), 
where \( h \) is the head entity, \( r \) the relation, and \( t \) the tail entity. Error detection involves determining if a given triple is correct or incorrect, with the output being a binary label.

A triple is \textbf{incorrect} if the head or tail entity does not align with the relation \cite{error}. For example, \textit{(Harvard University, is\_located\_in, New York)}. Conversely, a triple is \textbf{correct} if all components align appropriately. 


\subsection{Subgraph Definitions}

To analyze the context of a triple, we define two key concepts for each entity: \textit{Out\_Neighbor Subgraph} and \textit{In\_Neighbor Subgraph}.

\noindent\textbf{\textit{Out\_Neighbor Subgraph}}: The set of triples where the entity serves as the head. For an entity \( e \), the \textit{Out\_Neighbor Subgraph} is \( \{(e, r', t') \mid (e, r', t') \in \mathcal{G}\} \), where \( r' \) is outgoing relations from \( e \), and \( t' \) is the corresponding tail entity.

\noindent\textbf{\textit{In\_Neighbor Subgraph}}: The set of triples where the entity serves as the tail. For an entity \( e \), the \textit{In\_Neighbor Subgraph} is \( \{(h', r', e) \mid (h', r', e) \in \mathcal{G}\} \), where \( h' \) is the corresponding head entity, and \( r' \) represents incoming relations to \( e \).

Based on these concepts, for a given triple \( (h, r, t) \), we define the following subgraphs for both the head \( h \) and the tail \( t \):

\noindent\textbf{(a) \textit{Head\_Forward\_Subgraph}}: The \textit{Out\_Neighbor Subgraph} of the head entity \( h \), excluding the current triple \( (h, r, t) \). Formally:
\begin{align}
   \{(h, r', t') \mid (h, r', t') \in \mathcal{G}, (r', t') \neq (r, t)\}
\end{align}
   
\noindent\textbf{(b) \textit{Head\_Backward\_Subgraph}}: The \textit{In\_Neighbor Subgraph} of the head entity \( h \), capturing all incoming relations to \( h \). Formally:
\begin{align}
   \{(h', r', h) \mid (h', r', h) \in \mathcal{G} \}
\end{align}

\noindent\textbf{(c) \textit{Tail\_Forward\_Subgraph}}: The \textit{Out\_Neighbor Subgraph} of the tail entity \( t \), capturing all outgoing relations from \( t \). Formally:
\begin{align}
   \{(t, r', t') \mid (t, r', t') \in \mathcal{G}\}
\end{align}

\noindent\textbf{(d) \textit{Tail\_Backward\_Subgraph}}: The \textit{In\_Neighbor Subgraph} of the tail entity \( t \), excluding the current triple \( (h, r, t) \). Formally:
\begin{align}
   \{(h', r', t) \mid (h', r', t) \in \mathcal{G}, (h', r') \neq (h, r)\}
\end{align}


\subsection{Agent Construction}

We construct four agents based on the above subgraphs: 
\textit{Head\_Forward\_Agent}, \textit{Head\_Backward\_Agent}, \textit{Tail\_Forward\_Agent} and \textit{Tail\_Backward\_Agent}.
Each agent analyzes the corresponding subgraph for the triple \( (h, r, t) \), enabling a multi-angle evaluation of the triple by considering both head and tail entities' forward and backward contexts.


